% !TEX program = xelatex
% Compile with XeLaTeX. Requires the standard ctex/Fandol packages in TeX Live.
% 中文之后为英文；独立公式共用。Chinese text is followed by English.
\documentclass[11pt,a4paper]{article}

% --- Page margins and paragraphs (based on the supplied template) ---
\usepackage[margin=1in]{geometry}
\usepackage{parskip}
\usepackage[UTF8,scheme=plain,fontset=fandol]{ctex}
\usepackage{fontspec}
\setmainfont{Latin Modern Roman}

% --- Basic mathematics and typesetting ---
\usepackage{amsmath,amssymb,amsthm,mathtools}
\usepackage{xcolor}
\usepackage{longtable,booktabs,array,calc}
\usepackage{enumitem}
\usepackage{fancyhdr}
\usepackage{titlesec}
\usepackage{needspace}
\usepackage[hidelinks,unicode]{hyperref}
\definecolor{TranslationBlue}{HTML}{36546C}
\definecolor{NoteGray}{HTML}{505963}

% --- Definition and theorem environments, retained for further editing ---
\theoremstyle{definition}
\newtheorem{definition}{定义 / Definition}
\newtheorem{example}{例 / Example}
\theoremstyle{plain}
\newtheorem{theorem}{定理 / Theorem}
\newtheorem{proposition}{命题 / Proposition}
\newtheorem{lemma}{引理 / Lemma}

% --- Bilingual paragraph styles ---
\newenvironment{englishtranslation}
  {\par\begingroup\small\color{TranslationBlue}\setlength{\parskip}{3pt}}
  {\par\endgroup}
\newenvironment{editorialnote}
  {\par\begingroup\small\color{NoteGray}\setlength{\parskip}{4pt}}
  {\par\endgroup\medskip}
\newenvironment{bilingualpair}
  {\par\noindent\begin{minipage}{\linewidth}\setlength{\parskip}{4pt}\interdisplaylinepenalty=10000}
  {\end{minipage}\par\addvspace{4pt}}
\setlength{\parskip}{4pt plus 1pt}
\setlength{\parindent}{0pt}
\linespread{1.06}
\setlength{\emergencystretch}{2em}
\setlength{\abovedisplayskip}{8pt plus 2pt minus 2pt}
\setlength{\belowdisplayskip}{8pt plus 2pt minus 2pt}
\setlength{\jot}{4pt}
\allowdisplaybreaks[1]
\raggedbottom
\setcounter{secnumdepth}{0}
\setlist{topsep=4pt,itemsep=3pt,parsep=2pt,leftmargin=1.6em}
\providecommand{\tightlist}{\setlength{\itemsep}{0pt}\setlength{\parskip}{0pt}}
\providecommand{\passthrough}[1]{#1}

\titleformat{\section}{\large\bfseries}{}{0pt}{}
\titleformat{\subsection}{\normalsize\bfseries}{}{0pt}{}
\titleformat{\subsubsection}{\normalsize\bfseries}{}{0pt}{}
\titlespacing*{\section}{0pt}{16pt plus 3pt}{8pt}
\titlespacing*{\subsection}{0pt}{12pt plus 2pt}{6pt}
\titlespacing*{\subsubsection}{0pt}{10pt plus 2pt}{5pt}
\pagestyle{fancy}
\fancyhf{}
\fancyhead[L]{\small 数学物理方法 / Mathematical Methods for Physics}
\fancyhead[R]{\small Lecture 14}
\fancyfoot[C]{\small\thepage}
\setlength{\headheight}{15pt}
\renewcommand{\headrulewidth}{0.3pt}
\hypersetup{pdftitle={Mathematical Methods for Physics - Lecture 14}}

\begin{document}
\begin{center}
  {\LARGE\bfseries 数学物理方法 · 第 14 讲}\\[0.5em]
  {\Large\bfseries Mathematical Methods for Physics}\\[0.2em]
  {\large Lecture 14}\\[0.7em]
  {\normalsize \(\delta\) 函数、广义傅里叶变换及其性质}\\[0.2em]
  {\small\color{TranslationBlue} The Delta Function, Generalized Fourier Transforms, and Their Properties}
\end{center}
\vspace{0.5em}

\Needspace{12\baselineskip}

\section{\texorpdfstring{7.3 \(\delta\)
函数与广义傅里叶变换\\[2pt]{\normalfont\normalsize\color{TranslationBlue} 7.3
The Delta Function and Generalized Fourier
Transforms}}{7.3 \textbackslash delta 函数与广义傅里叶变换7.3 The Delta Function and Generalized Fourier Transforms}}

\Needspace{8\baselineskip}

\subsection{\texorpdfstring{1. \(\delta\)
函数与矩形脉冲近似\\[2pt]{\normalfont\small\color{TranslationBlue} The
Delta Function and Its Rectangular-Pulse
Approximation}}{1. \textbackslash delta 函数与矩形脉冲近似The Delta Function and Its Rectangular-Pulse Approximation}}

\begin{bilingualpair}

\(\delta\) 函数又叫单位脉冲函数。下式作直观描述：

\nopagebreak[4]\begin{englishtranslation}

The \(\delta\) function is also called the unit impulse. The original
notes give the following intuitive description.

\end{englishtranslation}

\[
\delta(x)=
\begin{cases}
0,&x\ne0,\\
\infty,&x=0,
\end{cases}
\qquad
\int_{-\infty}^{+\infty}\delta(x)\,\mathrm dx=1.
\]

\end{bilingualpair}

\begin{bilingualpair}

\textbf{定义 1}

\nopagebreak[4]\begin{englishtranslation}

\textbf{Definition 1}

\end{englishtranslation}

\[
\delta(x)=\lim_{\tau\to0}\delta_\tau(x),
\]

\end{bilingualpair}

\begin{bilingualpair}

其中

\nopagebreak[4]\begin{englishtranslation}

where

\end{englishtranslation}

\[
\delta_\tau(x)=
\begin{cases}
\dfrac1{\tau},&|x|\le\dfrac{\tau}{2},\\[4pt]
0,&|x|>\dfrac{\tau}{2}.
\end{cases}
\]

\end{bilingualpair}

\begin{bilingualpair}

\textbf{原图说明}：矩形脉冲以原点为中心，宽度为 \(\tau\)，高度为
\(1/\tau\)，面积为 1；当 \(\tau\to0\)
时，脉冲越来越窄、越来越高。\(\delta(x)\)
用原点处的竖直箭头表示，旁标``1''表示单位积分。

\nopagebreak[4]\begin{englishtranslation}

\textbf{Description of the original figure.} The rectangular pulse is
centered at the origin, has width \(\tau\) and height \(1/\tau\), and
has unit area. As \(\tau\to0\), it becomes narrower and taller. The
delta function is represented by a vertical arrow at the origin; the
label ``1'' denotes its unit integral.

\end{englishtranslation}

\end{bilingualpair}

\begin{bilingualpair}
Notes:
\par
\vspace*{3cm}
\end{bilingualpair}

\Needspace{8\baselineskip}

\subsection{\texorpdfstring{2. \(\delta\)
函数的泛函定义\\[2pt]{\normalfont\small\color{TranslationBlue} The Delta
Function as a
Functional}}{2. \textbackslash delta 函数的泛函定义The Delta Function as a Functional}}

\begin{bilingualpair}

\textbf{定义 2（数学定义）}

\nopagebreak[4]\begin{englishtranslation}

\textbf{Definition 2 (mathematical definition)}

\end{englishtranslation}

\end{bilingualpair}

\begin{bilingualpair}

\(\delta(t)\) 是一个``泛函''，满足

\nopagebreak[4]\begin{englishtranslation}

The delta function \(\delta(t)\) is a functional satisfying

\end{englishtranslation}

\[
\int_{-\infty}^{+\infty}
\delta(t-t_0)\varphi(t)\,\mathrm dt
=\varphi(t_0),
\qquad \forall\varphi\in C(\mathbb R).
\]

\end{bilingualpair}

\begin{bilingualpair}

其导数满足

\nopagebreak[4]\begin{englishtranslation}

Its derivative satisfies

\end{englishtranslation}

\[
\int_{-\infty}^{+\infty}
\delta'(t-t_0)\varphi(t)\,\mathrm dt
=-\varphi'(t_0),
\qquad \forall\varphi\in C^1(\mathbb R).
\]

\end{bilingualpair}

\begin{bilingualpair}

其 \(n\) 阶导数满足

\nopagebreak[4]\begin{englishtranslation}

Its \(n\)th derivative satisfies

\end{englishtranslation}

\[
\int_{-\infty}^{+\infty}
\delta^{(n)}(t-t_0)\varphi(t)\,\mathrm dt
=(-1)^n\varphi^{(n)}(t_0),
\qquad \forall\varphi\in C^n(\mathbb R).
\]

\end{bilingualpair}

\begin{editorialnote}

\end{editorialnote}

\begin{bilingualpair}
Notes:
\par
\vspace*{3cm}
\end{bilingualpair}

\Needspace{8\baselineskip}

\subsection{\texorpdfstring{3. \(\delta\)
函数的傅里叶变换与积分表示\\[2pt]{\normalfont\small\color{TranslationBlue} Fourier
Transform and Integral Representation of the Delta
Function}}{3. \textbackslash delta 函数的傅里叶变换与积分表示Fourier Transform and Integral Representation of the Delta Function}}

\begin{bilingualpair}

\textbf{例 1}：\(\delta(t)\) 的傅里叶变换称为广义傅里叶变换。

\nopagebreak[4]\begin{englishtranslation}

\textbf{Example 1.} The Fourier transform of \(\delta(t)\) is understood
as a generalized Fourier transform.

\end{englishtranslation}

\[
\mathcal F[\delta(t)]
=\int_{-\infty}^{+\infty}\delta(t)e^{-i\omega t}\,\mathrm dt
=\left.e^{-i\omega t}\right|_{t=0}
=1.
\]

\end{bilingualpair}

\begin{bilingualpair}

因此 \(\{\delta(t),1\}\) 构成一个傅里叶变换对。由反变换，

\nopagebreak[4]\begin{englishtranslation}

Thus \(\{\delta(t),1\}\) is a Fourier transform pair. By the inverse
transform,

\end{englishtranslation}

\[
\delta(t)=\frac1{2\pi}
\int_{-\infty}^{+\infty}1\cdot e^{i\omega t}\,\mathrm d\omega.
\]

\end{bilingualpair}

\begin{bilingualpair}

\textbf{定义 3（\(\delta(t)\) 的积分定义）}

\nopagebreak[4]\begin{englishtranslation}

\textbf{Definition 3 (integral representation of \(\delta(t)\))}

\end{englishtranslation}

\[
\delta(t)=\frac1{2\pi}
\int_{-\infty}^{+\infty}e^{i\omega t}\,\mathrm d\omega.
\]

\end{bilingualpair}

\begin{bilingualpair}

\textbf{注}

\nopagebreak[4]\begin{englishtranslation}

\textbf{Note}

\end{englishtranslation}

\[
\int_{-\infty}^{+\infty}e^{\pm i\omega t}\,\mathrm dt
=2\pi\delta(\omega).
\tag{*}
\]

\end{bilingualpair}

\begin{editorialnote}

\begin{bilingualpair}

\textbf{整理注}：上述积分等式在广义函数意义下理解。

\nopagebreak[4]\begin{englishtranslation}

\textbf{Editorial note.} These integral identities are interpreted in
the sense of distributions.

\end{englishtranslation}

\end{bilingualpair}

\end{editorialnote}

\begin{bilingualpair}

教材 P.340：常用广义傅里叶变换对。原稿首先列出第 10 对
\(\{\delta(t),1\}\)，并提示第 12、13、14、15 对。

\nopagebreak[4]\begin{englishtranslation}

Textbook, p.~340: common generalized Fourier transform pairs. The
original notes first list pair 10, \(\{\delta(t),1\}\), and then refer
to pairs 12, 13, 14, and 15.

\end{englishtranslation}

\end{bilingualpair}

\begin{bilingualpair}
Notes:
\par
\vspace*{3cm}
\end{bilingualpair}

\Needspace{8\baselineskip}

\subsection{\texorpdfstring{4.
复指数、余弦和正弦的变换\\[2pt]{\normalfont\small\color{TranslationBlue} Transforms
of Complex Exponentials, Cosines, and
Sines}}{4. 复指数、余弦和正弦的变换Transforms of Complex Exponentials, Cosines, and Sines}}

\begin{bilingualpair}

\textbf{例}：求证

\nopagebreak[4]\begin{englishtranslation}

\textbf{Example.} Prove that

\end{englishtranslation}

\[
e^{i\omega_0t}
\ \xleftrightarrow{\mathcal F}\
2\pi\delta(\omega-\omega_0).
\]

\end{bilingualpair}

\begin{bilingualpair}

由式 \((*)\)，

\nopagebreak[4]\begin{englishtranslation}

Using identity \((*)\),

\end{englishtranslation}

\[
\begin{aligned}
\mathcal F[e^{i\omega_0t}]
&=\int_{-\infty}^{+\infty}
e^{i\omega_0t}e^{-i\omega t}\,\mathrm dt\\
&=\int_{-\infty}^{+\infty}
e^{-i(\omega-\omega_0)t}\,\mathrm dt\\
&=2\pi\delta(\omega-\omega_0).
\end{aligned}
\]

\end{bilingualpair}

\begin{bilingualpair}

类似地，

\nopagebreak[4]\begin{englishtranslation}

Similarly,

\end{englishtranslation}

\[
\begin{aligned}
\mathcal F[\cos\omega_0t]
&=\mathcal F\!\left[
\frac{e^{i\omega_0t}+e^{-i\omega_0t}}2
\right]\\
&=\frac12\mathcal F[e^{i\omega_0t}]
+\frac12\mathcal F[e^{-i\omega_0t}]\\
&=\frac12\cdot2\pi\delta(\omega-\omega_0)
+\frac12\cdot2\pi\delta(\omega+\omega_0)\\
&=\pi\bigl[\delta(\omega+\omega_0)+\delta(\omega-\omega_0)\bigr].
\end{aligned}
\]

\end{bilingualpair}

\[
\mathcal F[\sin\omega_0t]
=i\pi\bigl[\delta(\omega+\omega_0)-\delta(\omega-\omega_0)\bigr].
\]

\begin{bilingualpair}
Notes:
\par
\vspace*{3cm}
\end{bilingualpair}

\Needspace{8\baselineskip}

\subsection{\texorpdfstring{5.
单位阶跃函数的变换\\[2pt]{\normalfont\small\color{TranslationBlue} Transform
of the Unit Step
Function}}{5. 单位阶跃函数的变换Transform of the Unit Step Function}}

\begin{bilingualpair}

\[
u(t)=
\begin{cases}
1,&t>0,\\
0,&t<0.
\end{cases}
\]

\end{bilingualpair}

\begin{bilingualpair}

可以验证

\nopagebreak[4]\begin{englishtranslation}

We can verify that

\end{englishtranslation}

\[
u(t)
\ \xleftrightarrow{\mathcal F}\
\frac1{i\omega}+\pi\delta(\omega).
\]

\end{bilingualpair}

\begin{bilingualpair}

从正变换形式上写出

\nopagebreak[4]\begin{englishtranslation}

The forward transform can formally be written as

\end{englishtranslation}

\[
\mathcal F[u(t)]
=\int_{-\infty}^{+\infty}u(t)e^{-i\omega t}\,\mathrm dt
=\int_0^{+\infty}e^{-i\omega t}\,\mathrm dt.
\]

\end{bilingualpair}

\begin{bilingualpair}

原稿接着用反变换验证：

\nopagebreak[4]\begin{englishtranslation}

The original notes then verify the result using the inverse transform.

\end{englishtranslation}

\[
\begin{aligned}
\mathcal F^{-1}\!\left[
\frac1{i\omega}+\pi\delta(\omega)
\right]
&=\frac1{2\pi}\int_{-\infty}^{+\infty}
\left[\frac1{i\omega}+\pi\delta(\omega)\right]
e^{i\omega t}\,\mathrm d\omega\\
&=\frac1{2\pi}\left[
\int_{-\infty}^{+\infty}
\frac{e^{i\omega t}}{i\omega}\,\mathrm d\omega
+\int_{-\infty}^{+\infty}
\pi\delta(\omega)e^{i\omega t}\,\mathrm d\omega
\right]\\
&=\frac1{2\pi}\left[
\int_{-\infty}^{+\infty}
\frac{e^{i\omega t}}{i\omega}\,\mathrm d\omega+\pi
\right].
\end{aligned}
\]

\end{bilingualpair}

\begin{bilingualpair}

其中

\nopagebreak[4]\begin{englishtranslation}

where

\end{englishtranslation}

\[
\begin{aligned}
\int_{-\infty}^{+\infty}
\frac{e^{i\omega t}}{i\omega}\,\mathrm d\omega
&=\int_{-\infty}^{+\infty}
\frac{\cos\omega t+i\sin\omega t}{i\omega}\,\mathrm d\omega\\
&=0+\int_{-\infty}^{+\infty}
\frac{\sin\omega t}{\omega}\,\mathrm d\omega\\
&=
\begin{cases}
\pi,&t>0,\\
-\pi,&t<0.
\end{cases}
\end{aligned}
\]

\end{bilingualpair}

\begin{bilingualpair}

\textbf{注}：由

\nopagebreak[4]\begin{englishtranslation}

\textbf{Note.} From

\end{englishtranslation}

\[
\int_{-\infty}^{+\infty}\frac{\sin x}{x}\,\mathrm dx=\pi
\]

\end{bilingualpair}

\begin{bilingualpair}

可得

\nopagebreak[4]\begin{englishtranslation}

we obtain

\end{englishtranslation}

\[
\int_{-\infty}^{+\infty}\frac{\sin kx}{x}\,\mathrm dx
=
\begin{cases}
\pi,&k>0,\\
-\pi,&k<0.
\end{cases}
\]

\end{bilingualpair}

\begin{bilingualpair}

代入上式，

\nopagebreak[4]\begin{englishtranslation}

Substituting into the preceding expression gives

\end{englishtranslation}

\[
\mathcal F^{-1}\!\left[
\frac1{i\omega}+\pi\delta(\omega)
\right]
=
\begin{cases}
\dfrac1{2\pi}(\pi+\pi)=1,&t>0,\\[4pt]
\dfrac1{2\pi}(-\pi+\pi)=0,&t<0,
\end{cases}
=u(t).
\]

\end{bilingualpair}

\begin{bilingualpair}
Notes:
\par
\vspace*{3cm}
\end{bilingualpair}

\Needspace{8\baselineskip}

\subsection{\texorpdfstring{6.
直流信号\\[2pt]{\normalfont\small\color{TranslationBlue} A Constant (DC)
Signal}}{6. 直流信号A Constant (DC) Signal}}

\begin{bilingualpair}

直流信号 \(E\) 的变换对为

\nopagebreak[4]\begin{englishtranslation}

The transform pair for a constant (DC) signal \(E\) is

\end{englishtranslation}

\[
E\ \xleftrightarrow{\mathcal F}\ 2\pi E\delta(\omega).
\]

\end{bilingualpair}

\[
\mathcal F[E]
=\int_{-\infty}^{+\infty}Ee^{-i\omega t}\,\mathrm dt
=2\pi E\delta(\omega).
\]

\begin{bilingualpair}
Notes:
\par
\vspace*{3cm}
\end{bilingualpair}

\Needspace{12\baselineskip}

\section{\texorpdfstring{二、傅里叶变换对的性质\\[2pt]{\normalfont\normalsize\color{TranslationBlue} Properties
of Fourier Transform
Pairs}}{二、傅里叶变换对的性质Properties of Fourier Transform Pairs}}

\begin{bilingualpair}

\emph{教材 P.190。}

\nopagebreak[4]\begin{englishtranslation}

\emph{textbook, p.~190.}

\end{englishtranslation}

\end{bilingualpair}

\begin{bilingualpair}

约定

\nopagebreak[4]\begin{englishtranslation}

We use the notation

\end{englishtranslation}

\[
f(t)\ \xleftrightarrow{\mathcal F}\ F(\omega),
\qquad
g(t)\ \xleftrightarrow{\mathcal F}\ G(\omega).
\]

\end{bilingualpair}

\begin{bilingualpair}

全文采用

\nopagebreak[4]\begin{englishtranslation}

Throughout these notes, the transform convention is

\end{englishtranslation}

\[
F(\omega)=\int_{-\infty}^{+\infty}f(t)e^{-i\omega t}\,\mathrm dt,
\qquad
f(t)=\frac1{2\pi}\int_{-\infty}^{+\infty}F(\omega)e^{i\omega t}\,\mathrm d\omega.
\]

\end{bilingualpair}

\Needspace{8\baselineskip}

\subsection{\texorpdfstring{1. 性质
1：线性性\\[2pt]{\normalfont\small\color{TranslationBlue} Property 1:
Linearity}}{1. 性质 1：线性性Property 1: Linearity}}

\begin{bilingualpair}

\[
\begin{aligned}
\mathcal F[af(t)+bg(t)]
&=a\mathcal F[f]+b\mathcal F[g]\\
&=aF(\omega)+bG(\omega).
\end{aligned}
\]

\end{bilingualpair}

\begin{bilingualpair}

或者

\nopagebreak[4]\begin{englishtranslation}

Alternatively,

\end{englishtranslation}

\[
\mathcal F^{-1}[aF(\omega)+bG(\omega)]
=a\mathcal F^{-1}[F]+b\mathcal F^{-1}[G].
\]

\end{bilingualpair}

\Needspace{8\baselineskip}

\subsubsection{\texorpdfstring{例
7.10\\[2pt]{\normalfont\small\color{TranslationBlue} Example
7.10}}{例 7.10Example 7.10}}

\[
\begin{aligned}
\mathcal F[\sin^2t]
&=\mathcal F\!\left[\frac{1-\cos2t}{2}\right]\\
&=\frac12\left\{\mathcal F[1]-\mathcal F[\cos2t]\right\}\\
&=\frac12\left\{
2\pi\delta(\omega)
-\pi[\delta(\omega+2)+\delta(\omega-2)]
\right\}.
\end{aligned}
\]

\Needspace{8\baselineskip}

\subsubsection{\texorpdfstring{例
7.11\\[2pt]{\normalfont\small\color{TranslationBlue} Example
7.11}}{例 7.11Example 7.11}}

\begin{bilingualpair}

设

\nopagebreak[4]\begin{englishtranslation}

Let

\end{englishtranslation}

\[
F(\omega)=\pi\delta(\omega+1)-\frac{i}{\omega+1},
\]

\end{bilingualpair}

\begin{bilingualpair}

求 \(\mathcal F^{-1}[F(\omega)]\)。

\nopagebreak[4]\begin{englishtranslation}

Find \(\mathcal F^{-1}[F(\omega)]\).

\end{englishtranslation}

\end{bilingualpair}

\begin{bilingualpair}

\textbf{解}

\nopagebreak[4]\begin{englishtranslation}

\textbf{Solution}

\end{englishtranslation}

\[
\begin{aligned}
\mathcal F^{-1}[F(\omega)]
&=\frac1{2\pi}\int_{-\infty}^{+\infty}
\left[\pi\delta(\omega+1)-\frac{i}{\omega+1}\right]
e^{i\omega t}\,\mathrm d\omega\\
&=\frac1{2\pi}\left[
\int_{-\infty}^{+\infty}
\pi\delta(\omega+1)e^{i\omega t}\,\mathrm d\omega
+\int_{-\infty}^{+\infty}
\frac{-ie^{i\omega t}}{\omega+1}\,\mathrm d\omega
\right]\\
&=\frac1{2\pi}\left[
\pi e^{-it}
+\int_{-\infty}^{+\infty}
\frac{-ie^{i\omega t}}{\omega+1}\,\mathrm d\omega
\right].
\end{aligned}
\]

\end{bilingualpair}

\begin{bilingualpair}

令 \(\omega+1=\xi\)，则

\nopagebreak[4]\begin{englishtranslation}

Let \(\omega+1=\xi\). Then

\end{englishtranslation}

\[
\begin{aligned}
\int_{-\infty}^{+\infty}
\frac{-ie^{i\omega t}}{\omega+1}\,\mathrm d\omega
&=\int_{-\infty}^{+\infty}
\frac{-ie^{i(\xi-1)t}}{\xi}\,\mathrm d\xi\\
&=e^{-it}\int_{-\infty}^{+\infty}
\frac{-ie^{i\xi t}}{\xi}\,\mathrm d\xi\\
&=e^{-it}
\begin{cases}
\pi,&t>0,\\
-\pi,&t<0.
\end{cases}
\end{aligned}
\]

\end{bilingualpair}

\begin{bilingualpair}

代入得

\nopagebreak[4]\begin{englishtranslation}

Substitution gives

\end{englishtranslation}

\[
\begin{aligned}
\mathcal F^{-1}[F(\omega)]
&=
\begin{cases}
\dfrac1{2\pi}(\pi e^{-it}+\pi e^{-it}),&t>0,\\[4pt]
\dfrac1{2\pi}(\pi e^{-it}-\pi e^{-it}),&t<0
\end{cases}\\
&=
\begin{cases}
e^{-it},&t>0,\\
0,&t<0
\end{cases}\\
&=u(t)e^{-it}.
\end{aligned}
\]

\end{bilingualpair}

\begin{bilingualpair}

即

\nopagebreak[4]\begin{englishtranslation}

That is,

\end{englishtranslation}

\[
u(t)e^{-it}
\ \xleftrightarrow{\mathcal F}\
\pi\delta(\omega+1)-\frac{i}{\omega+1}.
\]

\end{bilingualpair}

\begin{bilingualpair}

与

\nopagebreak[4]\begin{englishtranslation}

Compare this with

\end{englishtranslation}

\[
u(t)\ \xleftrightarrow{\mathcal F}\
\pi\delta(\omega)+\frac1{i\omega}
\]

\end{bilingualpair}

\begin{bilingualpair}

对比，频域发生了平移。

\nopagebreak[4]\begin{englishtranslation}

The comparison shows a shift in the frequency domain.

\end{englishtranslation}

\end{bilingualpair}

\begin{bilingualpair}
Notes:
\par
\vspace*{3cm}
\end{bilingualpair}

\Needspace{8\baselineskip}

\subsection{\texorpdfstring{2. 性质
2：平移性\\[2pt]{\normalfont\small\color{TranslationBlue} Property 2:
Shifting}}{2. 性质 2：平移性Property 2: Shifting}}

\begin{bilingualpair}

\[
f(t-t_0)
\ \xleftrightarrow{\mathcal F}\
e^{-it_0\omega}F(\omega).
\]

\end{bilingualpair}

\begin{bilingualpair}

或者

\nopagebreak[4]\begin{englishtranslation}

Alternatively,

\end{englishtranslation}

\[
e^{i\omega_0t}f(t)
\ \xleftrightarrow{\mathcal F}\
F(\omega-\omega_0).
\]

\end{bilingualpair}

\begin{bilingualpair}

由正、逆变换公式可得以上关系。

\nopagebreak[4]\begin{englishtranslation}

These relations follow from the forward and
inverse transform formulas.

\end{englishtranslation}

\end{bilingualpair}

\begin{bilingualpair}

\textbf{例}：由于

\nopagebreak[4]\begin{englishtranslation}

\textbf{Example.} Since

\end{englishtranslation}

\[
u(t)\ \xleftrightarrow{\mathcal F}\
\frac1{i\omega}+\pi\delta(\omega),
\]

\end{bilingualpair}

\begin{bilingualpair}

故

\nopagebreak[4]\begin{englishtranslation}

Hence,

\end{englishtranslation}

\[
u(t-c)\ \xleftrightarrow{\mathcal F}\
\left[\frac1{i\omega}+\pi\delta(\omega)\right]e^{-ic\omega}.
\]

\end{bilingualpair}

\begin{bilingualpair}

可进一步化简为

\nopagebreak[4]\begin{englishtranslation}

This can be simplified to

\end{englishtranslation}

\[
\frac{e^{-ic\omega}}{i\omega}+\pi\delta(\omega).
\]

\end{bilingualpair}

\begin{bilingualpair}

这里使用了 \(\delta\) 函数的性质

\nopagebreak[4]\begin{englishtranslation}

Here we use the delta-function identity

\end{englishtranslation}

\[
\delta(x-x_0)\varphi(x)
=\delta(x-x_0)\varphi(x_0).
\]

\end{bilingualpair}

\begin{bilingualpair}

另外，

\nopagebreak[4]\begin{englishtranslation}

In addition,

\end{englishtranslation}

\[
e^{-it}u(t)\ \xleftrightarrow{\mathcal F}\
\pi\delta(\omega+1)+\frac1{i(\omega+1)},
\]

\end{bilingualpair}

\begin{bilingualpair}

以及

\nopagebreak[4]\begin{englishtranslation}

and

\end{englishtranslation}

\[
\mathcal F^{-1}\!\left[
\pi\delta(\omega-1)+\frac1{i(\omega-1)}
\right]
=u(t)e^{it}.
\]

\end{bilingualpair}

\begin{bilingualpair}
Notes:
\par
\vspace*{3cm}
\end{bilingualpair}

\Needspace{8\baselineskip}

\subsection{\texorpdfstring{3. 性质
3：相似性\\[2pt]{\normalfont\small\color{TranslationBlue} Property 3:
Scaling}}{3. 性质 3：相似性Property 3: Scaling}}

\begin{bilingualpair}

若 \(f(t)\xleftrightarrow{\mathcal F}F(\omega)\)，则

\nopagebreak[4]\begin{englishtranslation}

If \(f(t)\xleftrightarrow{\mathcal F}F(\omega)\), then

\end{englishtranslation}

\[
f(at)\ \xleftrightarrow{\mathcal F}\
\frac1{|a|}F\!\left(\frac{\omega}{a}\right),
\qquad a\ne0.
\]

\end{bilingualpair}

\begin{bilingualpair}

并且

\nopagebreak[4]\begin{englishtranslation}

and

\end{englishtranslation}

\[
\frac1{|a|}f\!\left(\frac{t}{a}\right)
\ \xleftrightarrow{\mathcal F}\
F(a\omega).
\]

\end{bilingualpair}

\Needspace{8\baselineskip}

\subsubsection{\texorpdfstring{例：求
\(\mathcal F[u(5t-2)]\)\\[2pt]{\normalfont\small\color{TranslationBlue} Example:
Find
\(\mathcal F[u(5t-2)]\)}}{例：求 \textbackslash mathcal F{[}u(5t-2){]}Example: Find \textbackslash mathcal F{[}u(5t-2){]}}}

\begin{bilingualpair}

\textbf{解（定义法）}

\nopagebreak[4]\begin{englishtranslation}

\textbf{Solution (using the definition)}

\end{englishtranslation}

\[
\mathcal F[u(5t-2)]
=\int_{-\infty}^{+\infty}
u(5t-2)e^{-i\omega t}\,\mathrm dt.
\]

\end{bilingualpair}

\begin{bilingualpair}

令 \(5t-2=\tau\)，即 \(t=(\tau+2)/5\)，得

\nopagebreak[4]\begin{englishtranslation}

Let \(5t-2=\tau\), so that \(t=(\tau+2)/5\). We obtain

\end{englishtranslation}

\[
\begin{aligned}
\mathcal F[u(5t-2)]
&=\int_{-\infty}^{+\infty}
u(\tau)e^{-i\omega(\tau+2)/5}\frac15\,\mathrm d\tau\\
&=\frac15e^{-2i\omega/5}
\int_{-\infty}^{+\infty}
u(\tau)e^{-i(\omega/5)\tau}\,\mathrm d\tau\\
&=\frac15e^{-2i\omega/5}
\left[
\frac1{i(\omega/5)}+\pi\delta\!\left(\frac{\omega}{5}\right)
\right]\\
&=\frac{e^{-2i\omega/5}}{i\omega}
+\frac{\pi}{5}\delta\!\left(\frac{\omega}{5}\right).
\end{aligned}
\]

\end{bilingualpair}

\begin{editorialnote}

\begin{bilingualpair}

\textbf{注}：利用
\(\delta(\omega/5)=5\delta(\omega)\)，还可将第二项写成
\(\pi\delta(\omega)\)。

\nopagebreak[4]\begin{englishtranslation}

\textbf{note.} Using
\(\delta(\omega/5)=5\delta(\omega)\), the second term can also be
written as \(\pi\delta(\omega)\).

\end{englishtranslation}

\end{bilingualpair}

\end{editorialnote}

\begin{bilingualpair}
Notes:
\par
\vspace*{3cm}
\end{bilingualpair}

\Needspace{8\baselineskip}

\subsection{\texorpdfstring{4. 性质
4：对称性\\[2pt]{\normalfont\small\color{TranslationBlue} Property 4:
Duality}}{4. 性质 4：对称性Property 4: Duality}}

\begin{bilingualpair}

若

\nopagebreak[4]\begin{englishtranslation}

If

\end{englishtranslation}

\[
f(t)\ \xleftrightarrow{\mathcal F}\ F(\omega),
\]

\end{bilingualpair}

\begin{bilingualpair}

则

\nopagebreak[4]\begin{englishtranslation}

then

\end{englishtranslation}

\[
F(t)\ \xleftrightarrow{\mathcal F}\ 2\pi f(-\omega).
\]

\end{bilingualpair}

\begin{bilingualpair}

\textbf{证明}

\nopagebreak[4]\begin{englishtranslation}

\textbf{Proof}

\end{englishtranslation}

\[
F(\omega)
=\int_{-\infty}^{+\infty}f(t)e^{-i\omega t}\,\mathrm dt.
\]

\end{bilingualpair}

\begin{bilingualpair}

利用``定积分值与所用哑变量无关''，将上式中的 \(\omega\)、\(t\) 互换：

\nopagebreak[4]\begin{englishtranslation}

A definite integral does not depend on the name of its dummy variable.
Interchanging \(\omega\) and \(t\) in the notation gives

\end{englishtranslation}

\[
F(t)=\int_{-\infty}^{+\infty}f(\omega)e^{-it\omega}\,\mathrm d\omega.
\]

\end{bilingualpair}

\begin{bilingualpair}

令 \(\omega=-\xi\)，

\nopagebreak[4]\begin{englishtranslation}

Let \(\omega=-\xi\).

\end{englishtranslation}

\[
\begin{aligned}
F(t)
&=\int_{+\infty}^{-\infty}
f(-\xi)e^{it\xi}\,\mathrm d(-\xi)\\
&=\int_{-\infty}^{+\infty}
f(-\xi)e^{it\xi}\,\mathrm d\xi\\
&=\frac1{2\pi}\int_{-\infty}^{+\infty}
[2\pi f(-\omega)]e^{it\omega}\,\mathrm d\omega.
\end{aligned}
\]

\end{bilingualpair}

\begin{bilingualpair}

与反变换式比较，即得结论。

\nopagebreak[4]\begin{englishtranslation}

Comparison with the inverse transform formula proves the result.

\end{englishtranslation}

\end{bilingualpair}

\begin{bilingualpair}

\textbf{例}：由双边指数脉冲的变换对

\nopagebreak[4]\begin{englishtranslation}

\textbf{Example.} Start from the transform pair for a two-sided
exponential pulse,

\end{englishtranslation}

\[
e^{-\beta|t|}
\ \xleftrightarrow{\mathcal F}\
\frac{2\beta}{\omega^2+\beta^2},
\qquad \beta>0,
\]

\end{bilingualpair}

\begin{bilingualpair}

可得

\nopagebreak[4]\begin{englishtranslation}

we obtain

\end{englishtranslation}

\[
\frac{2\beta}{t^2+\beta^2}
\ \xleftrightarrow{\mathcal F}\
2\pi e^{-\beta|-\omega|}
=2\pi e^{-\beta|\omega|}.
\]

\end{bilingualpair}

\begin{bilingualpair}

因此

\nopagebreak[4]\begin{englishtranslation}

Therefore,

\end{englishtranslation}

\[
\frac1{t^2+\beta^2}
\ \xleftrightarrow{\mathcal F}\
\frac{\pi}{\beta}e^{-\beta|\omega|}.
\]

\end{bilingualpair}

\begin{bilingualpair}

\textbf{原图说明}：双边指数脉冲关于纵轴对称，在原点达到最大值，并向两侧衰减。

\nopagebreak[4]\begin{englishtranslation}

\textbf{Description of the original figure.} The two-sided exponential
pulse is symmetric about the vertical axis, reaches its maximum at the
origin, and decays in both directions.

\end{englishtranslation}

\end{bilingualpair}

\begin{editorialnote}

\begin{bilingualpair}

\textbf{注}：\(\beta>0\) 是所列衰减指数变换对的适用条件，此处补明。

\nopagebreak[4]\begin{englishtranslation}

\textbf{note.} The condition \(\beta>0\) has been made
explicit because it is required for the decaying-exponential transform
pair.

\end{englishtranslation}

\end{bilingualpair}

\end{editorialnote}

\begin{bilingualpair}
Notes:
\par
\vspace*{3cm}
\end{bilingualpair}

\Needspace{8\baselineskip}

\subsection{\texorpdfstring{5. 性质
5：微分性\\[2pt]{\normalfont\small\color{TranslationBlue} Property 5:
Differentiation}}{5. 性质 5：微分性Property 5: Differentiation}}

\begin{bilingualpair}

\textbf{（1）时域微分}

\nopagebreak[4]\begin{englishtranslation}

\textbf{(1) Differentiation in the time domain}

\end{englishtranslation}

\end{bilingualpair}

\begin{bilingualpair}

若 \(f(\pm\infty)=0\)，则

\nopagebreak[4]\begin{englishtranslation}

If \(f(\pm\infty)=0\), then

\end{englishtranslation}

\[
f'(t)\ \xleftrightarrow{\mathcal F}\
i\omega F(\omega).
\]

\end{bilingualpair}

\begin{bilingualpair}

\textbf{证明}：分部积分，

\nopagebreak[4]\begin{englishtranslation}

\textbf{Proof.} Integration by parts gives

\end{englishtranslation}

\[
\begin{aligned}
\int_{-\infty}^{+\infty}f'(t)e^{-i\omega t}\,\mathrm dt
&=\left.f(t)e^{-i\omega t}\right|_{-\infty}^{+\infty}
-\int_{-\infty}^{+\infty}(-i\omega)f(t)e^{-i\omega t}\,\mathrm dt\\
&=0+i\omega\int_{-\infty}^{+\infty}
f(t)e^{-i\omega t}\,\mathrm dt\\
&=i\omega F(\omega).
\end{aligned}
\]

\end{bilingualpair}

\begin{bilingualpair}

\textbf{（2）频域微分}

\nopagebreak[4]\begin{englishtranslation}

\textbf{(2) Differentiation in the frequency domain}

\end{englishtranslation}

\[
-itf(t)\ \xleftrightarrow{\mathcal F}\ F'(\omega).
\]

\end{bilingualpair}

\begin{bilingualpair}

\textbf{证明}

\nopagebreak[4]\begin{englishtranslation}

\textbf{Proof}

\end{englishtranslation}

\[
F(\omega)=\int_{-\infty}^{+\infty}
f(t)e^{-i\omega t}\,\mathrm dt,
\]

\end{bilingualpair}

\[
\begin{aligned}
F'(\omega)
&=\int_{-\infty}^{+\infty}
\frac{\partial}{\partial\omega}
[f(t)e^{-i\omega t}]\,\mathrm dt\\
&=\int_{-\infty}^{+\infty}
f(t)(-it)e^{-i\omega t}\,\mathrm dt.
\end{aligned}
\]

\begin{bilingualpair}

\textbf{推广}

\nopagebreak[4]\begin{englishtranslation}

\textbf{Generalization}

\end{englishtranslation}

\end{bilingualpair}

\begin{bilingualpair}

若

\nopagebreak[4]\begin{englishtranslation}

If

\end{englishtranslation}

\[
f(\pm\infty)=f'(\pm\infty)=\cdots
=f^{(n-1)}(\pm\infty)=0,
\]

\end{bilingualpair}

\begin{bilingualpair}

则

\nopagebreak[4]\begin{englishtranslation}

then

\end{englishtranslation}

\[
f^{(n)}(t)\ \xleftrightarrow{\mathcal F}\
(i\omega)^nF(\omega).
\]

\end{bilingualpair}

\begin{bilingualpair}

另一方面，

\nopagebreak[4]\begin{englishtranslation}

On the other hand,

\end{englishtranslation}

\[
(-it)^nf(t)\ \xleftrightarrow{\mathcal F}\
F^{(n)}(\omega).
\]

\end{bilingualpair}

\begin{bilingualpair}

由

\nopagebreak[4]\begin{englishtranslation}

From

\end{englishtranslation}

\[
\mathcal F[(-it)^nf(t)]=F^{(n)}(\omega)
\]

\end{bilingualpair}

\begin{bilingualpair}

可得

\nopagebreak[4]\begin{englishtranslation}

we obtain

\end{englishtranslation}

\[
\mathcal F[t^nf(t)]=i^nF^{(n)}(\omega).
\]

\end{bilingualpair}

\begin{bilingualpair}

特别地，

\nopagebreak[4]\begin{englishtranslation}

In particular,

\end{englishtranslation}

\[
\begin{aligned}
\mathcal F[t^n\cdot1]
&=i^n[2\pi\delta(\omega)]^{(n)}\\
&=2\pi i^n\delta^{(n)}(\omega).
\end{aligned}
\]

\end{bilingualpair}

\begin{bilingualpair}

即

\nopagebreak[4]\begin{englishtranslation}

That is,

\end{englishtranslation}

\[
t^n\ \xleftrightarrow{\mathcal F}\
2\pi i^n\delta^{(n)}(\omega).
\]

\end{bilingualpair}

\begin{bilingualpair}
Notes:
\par
\vspace*{3cm}
\end{bilingualpair}

\Needspace{8\baselineskip}

\subsection{\texorpdfstring{6.
广义傅里叶变换对举例\\[2pt]{\normalfont\small\color{TranslationBlue} Examples
of Generalized Fourier Transform
Pairs}}{6. 广义傅里叶变换对举例Examples of Generalized Fourier Transform Pairs}}

\begin{bilingualpair}

\emph{教材 P.340，第 16---32
对中的若干项。}

\nopagebreak[4]\begin{englishtranslation}

\emph{selected entries from pairs
16--32 on textbook p.~340.}

\end{englishtranslation}

\end{bilingualpair}

\begin{bilingualpair}

右栏的导数均对 \(\omega\) 求导。

\nopagebreak[4]\begin{englishtranslation}

All derivatives in the right-hand column are with respect to \(\omega\).

\end{englishtranslation}

\end{bilingualpair}

\Needspace{20\baselineskip}\begingroup\small\renewcommand{\arraystretch}{1.65}\everymath{\displaystyle}

\begin{longtable}[]{@{}
  >{\raggedright\arraybackslash}p{(\columnwidth - 4\tabcolsep) * \real{0.1600}}
  >{\raggedright\arraybackslash}p{(\columnwidth - 4\tabcolsep) * \real{0.2800}}
  >{\raggedright\arraybackslash}p{(\columnwidth - 4\tabcolsep) * \real{0.5600}}@{}}
\toprule\noalign{}
\begin{minipage}[b]{\linewidth}\raggedright
原编号 / Original no.
\end{minipage} & \begin{minipage}[b]{\linewidth}\raggedright
原函数 / Function
\end{minipage} & \begin{minipage}[b]{\linewidth}\raggedright
傅里叶变换 / Fourier transform
\end{minipage} \\
\midrule\noalign{}
\endhead
\bottomrule\noalign{}
\endlastfoot
16 & \(u(t-c)\) &
\(\left[\frac1{i\omega}+\pi\delta(\omega)\right]e^{-ic\omega}\) \\
17 & \(t\,u(t)\) &
\(i\frac{\mathrm d}{\mathrm d\omega}\left[\frac1{i\omega}+\pi\delta(\omega)\right]\) \\
18 & \(t^nu(t)\) &
\(i^n\frac{\mathrm d^n}{\mathrm d\omega^n}\left[\frac1{i\omega}+\pi\delta(\omega)\right]\) \\
21 & \(u(t)e^{iat}\) & \(\frac1{i(\omega-a)}+\pi\delta(\omega-a)\) \\
22 & \(u(t)e^{iat}t^n\) &
\(i^n\frac{\mathrm d^n}{\mathrm d\omega^n}\left[\frac1{i(\omega-a)}+\pi\delta(\omega-a)\right]\) \\
26 & \(\delta^{(n)}(t)\) & \((i\omega)^n\) \\
27 & \(\delta^{(n)}(t-c)\) & \((i\omega)^ne^{-ic\omega}\) \\
32 & \(t^ne^{iat}\) & \(2\pi i^n\delta^{(n)}(\omega-a)\) \\
\end{longtable}

\endgroup

\begin{bilingualpair}

第 21、32 对由频域平移 \(\omega\mapsto\omega-a\) 得到。

\nopagebreak[4]\begin{englishtranslation}

Pairs 21 and 32 follow from the frequency shift
\(\omega\mapsto\omega-a\).

\end{englishtranslation}

\end{bilingualpair}

\begin{bilingualpair}
Notes:
\par
\vspace*{3cm}
\end{bilingualpair}

\Needspace{8\baselineskip}

\subsection{\texorpdfstring{7. 性质
7：积分性\\[2pt]{\normalfont\small\color{TranslationBlue} Property 7:
Integration}}{7. 性质 7：积分性Property 7: Integration}}

\begin{bilingualpair}

\emph{性质编号沿用原稿，原稿此处未出现``性质 6''。}

\end{bilingualpair}

\begin{bilingualpair}

\textbf{（1）时域积分}

\nopagebreak[4]\begin{englishtranslation}

\textbf{(1) Integration in the time domain}

\end{englishtranslation}

\end{bilingualpair}

\begin{bilingualpair}

若

\nopagebreak[4]\begin{englishtranslation}

If

\end{englishtranslation}

\[
\int_{-\infty}^{+\infty}f(t)\,\mathrm dt=0,
\]

\end{bilingualpair}

\begin{bilingualpair}

则

\nopagebreak[4]\begin{englishtranslation}

then

\end{englishtranslation}

\[
\int_{-\infty}^{t}f(\tau)\,\mathrm d\tau
\ \xleftrightarrow{\mathcal F}\
\frac{F(\omega)}{i\omega}.
\]

\end{bilingualpair}

\begin{bilingualpair}

\textbf{（2）频域积分}

\nopagebreak[4]\begin{englishtranslation}

\textbf{(2) Integration in the frequency domain}

\end{englishtranslation}

\end{bilingualpair}

\begin{bilingualpair}
\[
\frac{f(t)}{-it}
\ \xleftrightarrow{\mathcal F}\
\int_{-\infty}^{\omega}F(\xi)\,\mathrm d\xi.
\]

\end{bilingualpair}

\begin{bilingualpair}

\textbf{证明（2）}

\nopagebreak[4]\begin{englishtranslation}

\textbf{Proof of (2)}

\end{englishtranslation}

\end{bilingualpair}

\begin{bilingualpair}

令

\nopagebreak[4]\begin{englishtranslation}

Let

\end{englishtranslation}

\[
G(\omega)=\int_{-\infty}^{\omega}F(\xi)\,\mathrm d\xi,
\]

\end{bilingualpair}

\begin{bilingualpair}

求导得 \(F(\omega)=G'(\omega)\)。由微分性，

\nopagebreak[4]\begin{englishtranslation}

Differentiation gives \(F(\omega)=G'(\omega)\). By the differentiation
property,

\end{englishtranslation}

\[
\begin{aligned}
\mathcal F^{-1}[F(\omega)]
&=\mathcal F^{-1}[G'(\omega)]\\
&=-it\,\mathcal F^{-1}[G(\omega)].
\end{aligned}
\]

\end{bilingualpair}

\begin{bilingualpair}

故

\nopagebreak[4]\begin{englishtranslation}

Hence,

\end{englishtranslation}

\[
\mathcal F^{-1}[G(\omega)]
=\frac{\mathcal F^{-1}[F(\omega)]}{-it}
=\frac{f(t)}{-it}.
\]

\end{bilingualpair}

\begin{bilingualpair}

即

\nopagebreak[4]\begin{englishtranslation}

That is,

\end{englishtranslation}

\[
\mathcal F^{-1}\!\left[
\int_{-\infty}^{\omega}F(\xi)\,\mathrm d\xi
\right]
=\frac{\mathcal F^{-1}[F]}{-it}.
\]

\end{bilingualpair}

\begin{bilingualpair}

\textbf{证明（1）}

\nopagebreak[4]\begin{englishtranslation}

\textbf{Proof of (1)}

\end{englishtranslation}

\end{bilingualpair}

\begin{bilingualpair}

令

\nopagebreak[4]\begin{englishtranslation}

Let

\end{englishtranslation}

\[
g(t)=\int_{-\infty}^{t}f(\tau)\,\mathrm d\tau,
\]

\end{bilingualpair}

\begin{bilingualpair}

则 \(f(t)=g'(t)\)。在 \(g(\pm\infty)=0\) 时，由微分性，

\nopagebreak[4]\begin{englishtranslation}

Then \(f(t)=g'(t)\). If \(g(\pm\infty)=0\), the differentiation property
gives

\end{englishtranslation}

\[
\mathcal F[f(t)]
=\mathcal F[g'(t)]
=i\omega\mathcal F[g(t)].
\]

\end{bilingualpair}

\begin{bilingualpair}

从而

\nopagebreak[4]\begin{englishtranslation}

It follows that

\end{englishtranslation}

\[
\mathcal F[g(t)]=\frac{F(\omega)}{i\omega}.
\]

\end{bilingualpair}

\begin{bilingualpair}

\textbf{注}：若

\nopagebreak[4]\begin{englishtranslation}

\textbf{Note.} If

\end{englishtranslation}

\[
\int_{-\infty}^{+\infty}f(t)\,\mathrm dt\ne0,
\]

\end{bilingualpair}

\begin{bilingualpair}

则

\nopagebreak[4]\begin{englishtranslation}

then

\end{englishtranslation}

\[
\mathcal F\!\left[
\int_{-\infty}^{t}f(\tau)\,\mathrm d\tau
\right]
=\frac{F(\omega)}{i\omega}
+\pi\delta(\omega)F(0).
\]

\end{bilingualpair}

\begin{bilingualpair}

用``卷积性''易得。

\nopagebreak[4]\begin{englishtranslation}

This follows readily from the convolution property.

\end{englishtranslation}

\end{bilingualpair}

\begin{bilingualpair}
Notes:
\par
\vspace*{3cm}
\end{bilingualpair}

\Needspace{8\baselineskip}

\Needspace{12\baselineskip}

\section{\texorpdfstring{三、作业\\[2pt]{\normalfont\normalsize\color{TranslationBlue} Assignments}}{三、作业记录Assignments}}

\Needspace{8\baselineskip}

\subsection{\texorpdfstring{11.22
作业\\[2pt]{\normalfont\small\color{TranslationBlue} Assignment for
November 22}}{11.22 作业Assignment for November 22}}

\begin{itemize}
\item
  \begin{bilingualpair}

  教材 P.340，第 8 对的推导。

  \nopagebreak[4]\begin{englishtranslation}

  Textbook, p.~340: derive transform pair 8.

  \end{englishtranslation}

  \end{bilingualpair}
\item
  \begin{bilingualpair}

  教材 P.203，第 1 题（双）、第 2 题、第 3 题（双）。

  \nopagebreak[4]\begin{englishtranslation}

  Textbook, p.~203: Exercise 1 (even-numbered parts), Exercise 2, and
  Exercise 3 (even-numbered parts).

  \end{englishtranslation}

  \end{bilingualpair}
\item
  \begin{bilingualpair}

  教材 P.203，第 7 题的（4）、（5）、（6）。

  \nopagebreak[4]\begin{englishtranslation}

  Textbook, p.~203: Exercise 7, parts (4), (5), and (6).

  \end{englishtranslation}

  \end{bilingualpair}
\end{itemize}

\Needspace{8\baselineskip}

\subsection{\texorpdfstring{11.29
作业\\[2pt]{\normalfont\small\color{TranslationBlue} Assignment for
November 29}}{11.29 作业Assignment for November 29}}

\begin{itemize}
\item
  \begin{bilingualpair}

  教材 P.203，第 4 题（双）。

  \nopagebreak[4]\begin{englishtranslation}

  Textbook, p.~203: Exercise 4 (even-numbered parts).

  \end{englishtranslation}

  \end{bilingualpair}
\item
  \begin{bilingualpair}

  教材 P.203，第 7 题的（1）、（2）、（3）。

  \nopagebreak[4]\begin{englishtranslation}

  Textbook, p.~203: Exercise 7, parts (1), (2), and (3).

  \end{englishtranslation}

  \end{bilingualpair}
\item
  \begin{bilingualpair}

  教材 P.203，第 6、8 题。

  \nopagebreak[4]\begin{englishtranslation}

  Textbook, p.~203: Exercises 6 and 8.

  \end{englishtranslation}

  \end{bilingualpair}
\item
  \nopagebreak[4]\begin{englishtranslation}

  B 套，第 7 题。

  \end{englishtranslation}

  \nopagebreak[4]\begin{englishtranslation}

  Set B, Exercise 7.

  \end{englishtranslation}
\end{itemize}

\begin{bilingualpair}
Notes:
\par
\vspace*{3cm}
\end{bilingualpair}

\end{document}
